% ============================================================
% DISCUSSION (P1 summary + takeaway; findings in figure order; novelty;
%             hypotheses restated; confounders; integration; limitations)
% ============================================================
\section{Discussion}
\label{sec:discussion}

% --- Paragraph 1: purpose, main finding, take-home ---
This study asked whether language models reproduce the racial and sex
disparities documented in human pain care when a patient's race and sex are
stated (H1) or implied by a name (H2). In Study~1, across four open-weight
models, six cases, eight identities and 2,880 scored responses, neither cue type
produced a difference that survived correction, and a change of identity moved
the opioid decision about as often as asking the same question twice. Study~2
extended the audit to \nConfigs{} configurations of \nModels{} proprietary and
open-weight models and to 50 clinician-written cases, and the answer held: no
difference survived the study-wide false discovery rate, every severity
difference was equivalent to zero within one point, and the few differences that
survived correction within a model favoured Black patients rather than
disadvantaging them. The take-home finding is that, in this constrained task,
current models assess pain and choose analgesia without a detectable disadvantage
to minority or female patients, while they differ among themselves, and across
their own reasoning and serving settings, by far more than they differ between
patients.

% --- Paragraph 2: explicit labels (Figure 2, left) ---
Under explicit labels (Figure~\ref{fig:differences}, left) the results support
H1. Every Holm-adjusted $p$-value was 1, the largest severity difference was a
tenth of a point, and the strong-opioid flip rate between identities was at or
below the same-prompt floor for every model. This replicates, on open-weight
models, the null that Young et al.\ reported for GPT-4 and Gemini
\cite{young2025racial}, and it does so with paired seeds and an explicit
sensitivity bound, which that study lacked. For severity the conclusion is stronger than non-significance: with the case as the unit, every difference was equivalent to zero within one point, half the smallest clinically important change \cite{farrar2001clinical}. For the opioid decision it is weaker: six cases leave intervals that cannot exclude differences of 20 percentage points, so the opioid null is an absence of evidence rather than evidence of absence. It is also consistent with the outcome-grounded triage audit of Jarrett
et al., in which race, language and insurance produced no harmful shift while
stigma-related attributes did \cite{jarrett2026outcome}.

% --- Paragraph 3: name cues (Figure 2, right; Appendix table); novelty; hypotheses restated ---
Under name cues (Figure~\ref{fig:differences}, right) the results do not support
H2. H2 predicted lower severity and fewer strong opioids for Black and Hispanic
names than for White male names; no model showed that direction. The novelty of
this condition lies in its two-name design, and the name-level analysis shows why
it mattered. gemma3's uniformly positive severity differences would, with one
name per group, have read as a model that rates every minority patient's pain
higher than a White man's. With two names, the shift resolved into a single
reference name, Greg Baker, that gemma3 rated lower than any other name (7.67
against 7.73--7.83), while the second White male name behaved like the rest and
a White female name was among the highest. A name carries more than race and sex
(it carries age cohort, class and regional associations \cite{bhetwal2026name}),
and a model's response to any one name is an idiosyncrasy until a second name of
the same group reproduces it. That is the same lesson Tonneau et al.\ drew at
scale, that cues for the same group are not interchangeable
\cite{tonneau2026cues}, and that Vohra and Ravikiran drew for audit instruments
\cite{vohra2026audit}: here the instrument that would have manufactured a
finding was the choice of a single reference name. Restating the hypotheses in
the light of the data: H1 was supported within the stated sensitivity, and H2 was
not supported in direction or magnitude.

% --- Paragraph 3b: Study 2 (Figures 4 and 5; Tables 7-9) ---
Study~2 addressed the third purpose, whether the answer depends on the cases or
the models chosen, and in direction it did not. On clinician-written cases, for
proprietary models from three developers and open-weight models from 3 billion
parameters to mixture-of-experts models of several hundred billion, with reasoning on and off, the human disparity did not appear. What
appeared instead, under explicit labels, was a small and consistent shift the
other way (Figure~\ref{fig:panel}): the Black minus White severity contrast was
positive in \nBlackSevPosExplicit{} of \nPooledExplicit{} configurations, and
every difference that survived correction within a model rated a Black patient's
pain higher, or offered an opioid more often, than the White male reference. The
shift was weaker under name cues (Figure~\ref{fig:panelnames}), as would be
expected if models respond to the explicit word rather than to anything a name
conveys, and it is consistent with alignment training that teaches models to
respond with care to a stated protected attribute \cite{ouyang2022training}. A
shift of a few hundredths of a point is clinically negligible, but it shows that
an explicit label does not leave the output untouched. Two other results carry
more practical weight than any identity effect. First, the models disagreed with
one another, and with the Q-Pain reference answer, far more than they
distinguished between patients: the same cases drew the high dose from llama3.2
in nearly every response and from gemma3 in almost none, and full agreement with
the clinical reference ranged from \rangeFullMatch{}. Second, the same weights
behaved differently under different settings: reasoning moved severity by up to
half a point, upward in some models and downward in others, and where it changed
opioid decisions it made opioids or high doses less frequent, and Gemma~3 served at 16-bit precision
through an API rather than at 4-bit locally recommended the high dose a quarter
more often while rating severity the same, a contrast that cannot separate
precision from the serving software. For a deployment, then, the model, its
reasoning setting and its serving stack determine what a patient is offered more
than the patient's race or sex does, and an audit of one configuration does not
certify another.

% --- Paragraph 4: confounders and alternative explanations (Figure 3) ---
Several alternative explanations were considered and addressed by design or
measurement. First, the pilot showed that when a vignette states the patient's
own rating, models copy it (60--100\% of responses), so an audit built on such
vignettes would find no disparity because no assessment takes place; the main
conditions therefore described observed behaviour only. Second, sampling
instability could either manufacture or hide a difference
\cite{atil2024nondeterminism, song2025good}; pairing every comparison on the
seed removed most of it, and reporting the same-prompt floor
(Figure~\ref{fig:noise}) and the matched different-seed comparison showed the residual counterfactual change close to the models' own noise, as Bellibatlu et al.\ recommend
\cite{bellibatlu2026instability}. gemma3 was the one model whose severity rating
changed between identities slightly more often than between seeds, in both
conditions, and in the explicit condition those changes ran in both directions,
which is the signature of a model that attends to the identity token without a
consistent valence. Third, because all seven identities are compared with the
same reference responses, the seven differences move together whenever a
reference response happens to flip; the uniformly negative opioid differences for
ministral-3 and llama3.2 under name cues are exactly what that construction
produces and were not treated as evidence. Fourth, the JSON constraint and the
single-sentence rationale limit the model's expressed reasoning, which may
suppress a bias that free text would reveal; the constraint was accepted
because it makes every response scorable without human judgement.

% --- Paragraph 5: fit with other research; translation ---
These findings fit a pattern across model generations. Early generative models
showed intersectional treatment gaps in pain management \cite{loge2021qpain},
2023 commercial models reproduced race-based misconceptions and stereotyped
vignettes \cite{omiye2023race, zack2024gpt4}, and the most recent proprietary
and now open-weight models show no measurable disparity in direct pain
assessment \cite{young2025racial}. The most economical reading is that
instruction tuning and alignment \cite{ouyang2022training} have removed the most
direct expression of this bias, at least when the task is short, the cues are
clean and the output is constrained. The same open-weight families are not bias
free elsewhere: they show name bias in fill-in-the-blank tasks
\cite{bhetwal2026name}, language-triggered bias in recruitment
\cite{kitahara2026linguistic}, and sensitivity to stigmatising language and
social attributes in triage \cite{huang2026artificial, jarrett2026outcome}, so
the null reported here is a statement about a task, not about the models. That
reading translates into practice in two ways. For deployment, it means that the
first-order risk of an open-weight model reproducing the human pain disparity is
not supported for these models in this task, while saying nothing about longer
notes, dialect \cite{hofmann2024dialect}, insurance status or free-text
reasoning. For auditing, it means that counterfactual audits of clinical LLMs
should adopt three practices that this study found decisive: several names per
group, comparisons paired on the sampling seed, and a reported noise floor and
minimum detectable effect, without which a difference of a tenth of a point
cannot be distinguished from a model's ordinary variability.

% --- Limitations paragraph ---
\subsection{Limitations}
\paragraph{Cases.} The six authored cases were drafted with a language model and
edited by the first author without clinician review, and six cases cannot
represent the range of pain presentations; the sickle cell case carries a
racial association in United States epidemiology that is independent of the
manipulation. Q-Pain's 50 cases were designed by one clinical expert, and only a
random subset was reviewed by other physicians. In all 50 the reference answer
is to offer the opioid at the low dose, so the set cannot show whether a model
withholds an opioid when withholding is correct, and a high-dose answer departs
from the reference rather than marking better care. Q-Pain has been public since
2021 and may be in the training data of the models audited; because the audit
fills its identity placeholders, any exposure applies to every identity arm of a
case alike, but a memorised answer could mute a model's response to the patient.
In both case sets the clinical case determined most of each model's severity
rating (86--100\% of variance in Study~1), which leaves little room for an
identity effect on severity; llama3.2's severity null in particular is
uninformative, since it returned one rating per case. Each case had a single
age, so age was not tested.

\paragraph{Identity cues.} Explicit labels and first names are only two of the
ways identity enters a clinical note; dialect, address, insurance and social
history were not tested. The authored Hispanic and Asian names were not
validated for perceived race, and Q-Pain's names, although drawn from
demographic name data \cite{tzioumis2018demographic}, were not validated for how
models perceive them. Names also signal age cohort, class and region, so a
name-level difference need not be racial. In the Q-Pain name condition the
template's pronouns carried sex alongside the name, so that condition removes
the explicit race word but not an explicit sex signal. The racial categories are
those of United States epidemiology, and the audit was English only.

\paragraph{Models and settings.} Every configuration ran at one temperature
(0.7), with one prompt wording and one physician role, as a single request with
no system prompt and no record context; deployed systems differ in all of these.
Each API model was reached through one provider, and another provider of the
same open weights may serve them at a different precision or with different
defaults; proprietary models' serving precision is not disclosed, nor is
Beam's on Ollama Cloud. Claude Haiku 4.5 and Beam could not be seeded, so their
pairs share a replicate index but not a seed, and Beam's reasoning could not be
switched off, so it has no reasoning contrast. On
the authored cases the local models answered in free JSON while the API models
answered under a schema, and the local reasoning configurations were run on the
authored cases only. The precision contrast compares a 4-bit copy served by Ollama with a 16-bit copy
served by DeepInfra, so precision is confounded with the serving software and its
prompt template. Models served through an API can change without notice;
the time of every request is recorded with its response.

\paragraph{Outcomes and inference.} The outcomes omit timing, route and
non-pharmacological management, which Young et al.\ examined
\cite{young2025racial}. In Study~1, thirty to sixty pairs per comparison from six
cases bound the detectable effect at roughly 0.1--0.2 severity points and 4--14
percentage points of strong-opioid use and leave the opioid estimates imprecise;
in Study~2 six cases cannot support a case-level test that survives correction.
The 50 Q-Pain cases narrow that limit without removing it, so disparities
smaller than the reported intervals remain possible. Neither study was
preregistered: H1 and H2 were stated after the Study~1 data were collected, and
Study~2 had no directional hypothesis, so the analyses are exploratory.
